% =============================================================================
% Chapter 15: Query Processing - Λυμένες Ασκήσεις
% Data Mining Study Guide - NTUA
% =============================================================================

\chapter{Query Processing - Λυμένες Ασκήσεις}
\label{ch:query-processing-exercises}

Αυτό το κεφάλαιο περιέχει λυμένες ασκήσεις για επεξεργασία ερωτημάτων και
βελτιστοποίηση.

% -----------------------------------------------------------------------------
% Exercise 4.1
% -----------------------------------------------------------------------------
\section{Άσκηση 4.1: Pipelining vs Materialization}

\begin{formulabox}[Εκφώνηση]
Για το query: $\pi_{name}(\sigma_{salary>50000}(Employee \bowtie Department))$

Ταξινομήστε τους τελεστές σε pipelined ή materialized και εξηγήστε γιατί.
\end{formulabox}

\subsection*{Λύση}

\textbf{Ορισμοί:}
\begin{itemize}
    \item \textbf{Pipelining:} Τα αποτελέσματα περνάνε αμέσως στον επόμενο τελεστή χωρίς αποθήκευση
    \item \textbf{Materialization:} Τα αποτελέσματα αποθηκεύονται προσωρινά πριν περάσουν στον επόμενο τελεστή
\end{itemize}

\begin{center}
\begin{tabular}{lll}
\toprule
\textbf{Τελεστής} & \textbf{Τύπος} & \textbf{Λόγος} \\
\midrule
$\bowtie$ (Join) & Materialized & Χρειάζεται πλήρη δεδομένα για hash/sort \\
$\sigma$ (Selection) & Pipelined & Φιλτράρει γραμμή-γραμμή \\
$\pi$ (Projection) & Pipelined & Επεξεργάζεται γραμμή-γραμμή \\
\bottomrule
\end{tabular}
\end{center}

\begin{infobox}[Pipeline Breakers]
Οι τελεστές που \textbf{δεν} μπορούν να είναι pipelined ονομάζονται «pipeline breakers»:
\begin{itemize}
    \item Sort
    \item Hash Join (build phase)
    \item Group By / Aggregation
    \item Duplicate elimination (DISTINCT)
\end{itemize}
\end{infobox}

% -----------------------------------------------------------------------------
% Exercise 4.2
% -----------------------------------------------------------------------------
\section{Άσκηση 4.2: Επιλογή Join Algorithm}

\begin{formulabox}[Εκφώνηση]
Δίνονται δύο πίνακες:
\begin{itemize}
    \item $R$: 10,000 tuples, 100 tuples/page $\Rightarrow$ 100 pages
    \item $S$: 1,000 tuples, 100 tuples/page $\Rightarrow$ 10 pages
    \item Buffer: 12 pages
\end{itemize}

Ποιος αλγόριθμος join είναι καλύτερος; Υπολογίστε το I/O cost.
\end{formulabox}

\subsection*{Λύση}

\textbf{Nested Loop Join:}

$$Cost_{NLJ} = B_R + B_R \cdot B_S = 100 + 100 \cdot 10 = 1100 \text{ I/Os}$$

(με το μικρότερο ως outer: $10 + 10 \cdot 100 = 1010$ I/Os)

\textbf{Block Nested Loop Join:}

Με buffer $B = 12$ pages, κρατάμε $B-2 = 10$ pages για τον outer.

$$Cost_{BNLJ} = B_R + \lceil B_R / (B-2) \rceil \cdot B_S$$

Με $S$ ως outer (μικρότερο):
$$Cost = 10 + \lceil 10/10 \rceil \cdot 100 = 10 + 100 = 110 \text{ I/Os}$$

\textbf{Hash Join:}

$$Cost_{Hash} = 3 \cdot (B_R + B_S) = 3 \cdot 110 = 330 \text{ I/Os}$$

\textbf{Sort-Merge Join:}

$$Cost_{SM} = Sort(R) + Sort(S) + B_R + B_S$$
$$= 2 \cdot 100 \cdot \lceil \log_{11}(10) \rceil + 2 \cdot 10 \cdot \lceil \log_{11}(1) \rceil + 110$$
$$\approx 200 \cdot 1 + 20 \cdot 1 + 110 = 330 \text{ I/Os}$$

\begin{center}
\begin{tabular}{lr}
\toprule
\textbf{Algorithm} & \textbf{I/O Cost} \\
\midrule
Simple NLJ (R outer) & 1,100 \\
Simple NLJ (S outer) & 1,010 \\
Block NLJ (S outer) & \textbf{110} \\
Hash Join & 330 \\
Sort-Merge & $\sim$330 \\
\bottomrule
\end{tabular}
\end{center}

\textbf{Τελική Απάντηση:} Block Nested Loop Join με S ως outer (110 I/Os).

% -----------------------------------------------------------------------------
% Exercise 4.3
% -----------------------------------------------------------------------------
\section{Άσκηση 4.3: Query Optimization (Push Selection)}

\begin{formulabox}[Εκφώνηση]
Βελτιστοποιήστε το query tree για:

$$\pi_{ename}(\sigma_{dname='Sales'}(Employee \bowtie Department))$$
\end{formulabox}

\subsection*{Λύση}

\textbf{Αρχικό Query Tree:}

\begin{center}
\begin{tikzpicture}[
    node distance=1.2cm,
    op/.style={circle, draw=secondary, fill=box-blue, minimum size=0.8cm},
    rel/.style={rectangle, draw=primary, fill=box-green, rounded corners, minimum width=1.5cm}]

    \node[op] (proj) {$\pi$};
    \node[op, below=of proj] (sel) {$\sigma$};
    \node[op, below=of sel] (join) {$\bowtie$};
    \node[rel, below left=1cm and 0.5cm of join] (emp) {Employee};
    \node[rel, below right=1cm and 0.5cm of join] (dept) {Department};

    \draw[arrow] (sel) -- (proj);
    \draw[arrow] (join) -- (sel);
    \draw[arrow] (emp) -- (join);
    \draw[arrow] (dept) -- (join);

    \node[right=0.3cm of proj, font=\small] {$name$};
    \node[right=0.3cm of sel, font=\small] {$dname='Sales'$};
\end{tikzpicture}
\end{center}

\textbf{Βελτιστοποιημένο Query Tree (Push Selection Down):}

\begin{center}
\begin{tikzpicture}[
    node distance=1.2cm,
    op/.style={circle, draw=secondary, fill=box-blue, minimum size=0.8cm},
    rel/.style={rectangle, draw=primary, fill=box-green, rounded corners, minimum width=1.5cm}]

    \node[op] (proj) {$\pi$};
    \node[op, below=of proj] (join) {$\bowtie$};
    \node[rel, below left=1cm and 0.5cm of join] (emp) {Employee};
    \node[op, below right=1cm and 0.5cm of join] (sel) {$\sigma$};
    \node[rel, below=of sel] (dept) {Department};

    \draw[arrow] (join) -- (proj);
    \draw[arrow] (emp) -- (join);
    \draw[arrow] (sel) -- (join);
    \draw[arrow] (dept) -- (sel);

    \node[right=0.3cm of proj, font=\small] {$name$};
    \node[right=0.3cm of sel, font=\small] {$dname='Sales'$};
\end{tikzpicture}
\end{center}

\begin{infobox}[Κανόνες Βελτιστοποίησης]
\begin{enumerate}
    \item \textbf{Push selections down}: Μειώνει το μέγεθος ενδιάμεσων αποτελεσμάτων
    \item \textbf{Push projections down}: Μειώνει το πλάτος tuples
    \item \textbf{Join ordering}: Μικρότερα ενδιάμεσα αποτελέσματα πρώτα
\end{enumerate}
\end{infobox}

% -----------------------------------------------------------------------------
% Exercise 4.4
% -----------------------------------------------------------------------------
\section{Άσκηση 4.4: Cost Comparison}

\begin{formulabox}[Εκφώνηση]
Για τον πίνακα $R$ με 10,000 tuples και index στο attribute $A$:

Συγκρίνετε το cost για: $\sigma_{A=5}(R)$

\begin{enumerate}
    \item Table scan
    \item B+-tree index (height = 3, clustering factor = 1.0)
    \item Hash index
\end{enumerate}
\end{formulabox}

\subsection*{Λύση}

Έστω $B_R = 100$ pages, selectivity = 1\% (100 tuples match).

\textbf{(a) Table Scan:}
$$Cost = B_R = 100 \text{ page I/Os}$$

\textbf{(b) B+-tree Index (clustered):}
$$Cost = height + \text{matching pages} = 3 + 1 = 4 \text{ I/Os}$$
(Clustered = τα matching tuples είναι συνεχόμενα)

\textbf{(c) Hash Index:}
$$Cost = 1.2 \text{ (average bucket access)} = \approx 2 \text{ I/Os}$$

\begin{center}
\begin{tabular}{lr}
\toprule
\textbf{Method} & \textbf{I/O Cost} \\
\midrule
Table Scan & 100 \\
B+-tree (clustered) & 4 \\
Hash Index & $\sim$2 \\
\bottomrule
\end{tabular}
\end{center}

\begin{warningbox}[Unclustered Index]
Αν ο B+-tree index ήταν unclustered, το cost θα ήταν:
$$Cost = 3 + 100 = 103 \text{ I/Os (χειρότερο από scan!)}$$
\end{warningbox}
